\documentclass{beamer}
\usepackage{amssymb}
\usepackage{amsmath}

\begin{document}

  \begin{frame}
    \frametitle{Definiciones basicas}
    \textit{Definición}: Un contexto formal es
    una terna $C = (O, P, R)$ donde $O$ es un
    conjunto de objetos, $P$ es un conjunto
    de atributos y $R \subseteq O \times P$
    es una relación binaria.

    \bigskip
    \pause

    \textit{Definición}: Dado un conjunto
    $A \subseteq P$, definimos:

    \pause
    \begin{itemize}[<+->]
      \item $A^\perp = \{x \in O : \forall a \in A,
        (x,a) \in R\}$ como la extensión de $A$.
      \item $A^{\perp\dagger} = (A^\perp)^\dagger$
        como la clausura intencional de $A$.
      \item Decimos que $A$ es un \textit{intento}
        si $A = A^{\perp\dagger}$.
    \end{itemize}
  \end{frame}

  \begin{frame}
    \frametitle{Implicaciones}
    \textit{Definición}: Una implicación sobre
    $P$ es una expresión $A \Rightarrow B$, con
    $A, B \subseteq P$.

    \bigskip
    \pause

    \textit{Teorema}: El contexto $C$ satisface
    $A \Rightarrow B$ (escrito $C \models A \Rightarrow B$)
    si y sólo si:

    \pause
    \begin{itemize}[<+->]
      \item $A^\perp \subseteq B^\perp$.
      \item $B \subseteq A^{\perp\dagger}$.
      \item La implicación $A \Rightarrow A^{\perp\dagger}$
        tiene toda la información de $A$.
    \end{itemize}
  \end{frame}

  \begin{frame}
    \frametitle{Pseudo-intentos}
    \textit{Definición}: Un conjunto $P \subseteq P$
    se llama \textit{pseudo-intento} si cumple:

    \bigskip
    \pause

    \begin{itemize}[<+->]
      \item $P \neq P^{\perp\dagger}$
        ($P$ no es cerrado).
      \item Para todo pseudo-intento $Q \subset P$,
        se cumple que $Q^{\perp\dagger} \subseteq P$.
    \end{itemize}

    \bigskip
    \pause

    \textit{Observación}: Son premisas mínimas que
    generan clausuras sin redundancia.
  \end{frame}

  \begin{frame}
    \frametitle{Base de Guigues-Duquenne}
    \textit{Definición}: La \textit{Base de GD} de un
    contexto $C$ es el conjunto:

    \[
      \mathcal{B} = \{P \Rightarrow P^{\perp\dagger} :
      P \text{ es pseudo-intento}\}
    \]

    \bigskip
    \pause

    \textit{Teorema}: La base $\mathcal{B}$ cumple:

    \pause
    \begin{itemize}[<+->]
      \item \textit{Validez}: Toda implicación en
        $\mathcal{B}$ se cumple en $C$.
      \item \textit{Completitud}: Toda regla válida
        en $C$ se deduce de $\mathcal{B}$.
      \item \textit{Minimalidad}: Es la base con el
        menor número de reglas posible.
    \end{itemize}
  \end{frame}

  \begin{frame}
    \frametitle{Ejemplo: contexto}
    Sea $C = (O, P, R)$ con objetos $O = \{1, 2, 3\}$
    y atributos $P = \{a, b, c, d\}$:

    \bigskip
    \pause

    \begin{center}
      \begin{tabular}{|c|c|c|c|c|}
        \hline
        \textbf{Objeto} & \textbf{a} & \textbf{b} &
        \textbf{c} & \textbf{d} \\ \hline
        1 & $\times$ & $\times$ & & \\ \hline
        2 & & $\times$ & $\times$ & \\ \hline
        3 & $\times$ & $\times$ & $\times$ & $\times$ \\ \hline
      \end{tabular}
    \end{center}

    \pause
    \begin{itemize}[<+->]
      \item Objeto 1: $\{a, b\}$.
      \item Objeto 2: $\{b, c\}$.
      \item Objeto 3: $\{a, b, c, d\}$.
    \end{itemize}
  \end{frame}

  \begin{frame}
    \frametitle{Ejemplo: calculos}
    Calculamos clausuras para hallar pseudo-intentos:

    \bigskip
    \pause

    \begin{itemize}[<+->]
      \item \textbf{Vacio $\emptyset$}: \\
        $\emptyset^{\perp\dagger} = \{a,b\} \cap \{b,c\}
        \cap \{a,b,c,d\} = \{b\}$. \\
        Como $\emptyset \neq \{b\}$, es pseudo-intento. \\
        $$\text{Regla 1: } \emptyset \Rightarrow \{b\}$$

      \item \textbf{Conjunto $\{b, d\}$}: \\
        Contiene a $\emptyset^{\perp\dagger} = \{b\}$. \\
        $\text{Clausura} = \{a, b, c, d\}$. \\
        Como $\{b,d\} \neq \{a,b,c,d\}$, es pseudo-intento. \\
        $$\text{Regla 2: } \{b, d\} \Rightarrow \{a, c, d\}$$

      \item \textbf{Conjunto $\{a, b, c\}$}: \\
        Contiene a $\{b\}$. \\
        $\text{Clausura} = \{a, b, c, d\}$. \\
        Es pseudo-intento. \\
        $$\text{Regla 3: } \{a, b, c\} \Rightarrow \{d\}$$
    \end{itemize}
  \end{frame}

  \begin{frame}
    \frametitle{Ejemplo: resultado}
    La base obtenida es:

    \[
      \mathcal{B} = \{
        \emptyset \Rightarrow \{b\}, \,
        \{b, d\} \Rightarrow \{a, c, d\}, \,
        \{a, b, c\} \Rightarrow \{d\}
      \}
    \]

    \bigskip
    \pause

    \textit{Significado}:

    \pause
    \begin{itemize}[<+->]
      \item $\emptyset \Rightarrow \{b\}$: Todos los objetos
        tienen el atributo $b$.
      \item $\{b, d\} \Rightarrow \{a, c, d\}$: Tener $b$ y $d$
        implica tener todos los demas.
      \item $\{a, b, c\} \Rightarrow \{d\}$: Tener $a,b,c$
        implica tener $d$.
      \item Resume todas las reglas sin repetirlas.
    \end{itemize}
  \end{frame}

\end{document}
